
\documentclass[11pt]{article}
\usepackage[margin=1in]{geometry}
\usepackage{amsmath, amssymb, amsthm}
\usepackage[T1]{fontenc}
\usepackage[utf8]{inputenc}
\usepackage{lmodern}
\usepackage{enumitem}
\usepackage{hyperref}

\title{The Parasite Machine:\\Toward a Relational Logic Without Substance}
\author{}
\date{}

\newtheorem{definition}{Definition}
\newtheorem{proposition}{Proposition}

\begin{document}
\maketitle

\section*{Abstract}
We introduce the \emph{Parasite Machine}, a minimal relational operator inspired by Michel Serres, capable of generating structure without presupposing substance, conservation, or fixed ontology. The machine operates on relations via a third term that may act as noise or metastability. We formalize its dynamics, define its transformation rules, and situate it within dynamical systems and relational logic.

\section{Motivation}

Classical modeling frameworks often assume:
\begin{itemize}
\item conserved quantities,
\item stable entities,
\item fixed ontological primitives.
\end{itemize}

We seek a minimal alternative:
\begin{quote}
A system defined only by relations, transformations, and emergent structure.
\end{quote}

\section{The Parasite Structure}

\begin{definition}[Parasite Machine]
A Parasite Machine is a tuple
\[
P = (A, B, C, d)
\]
where:
\begin{itemize}
\item $A \rightarrow B$ is a directed relation,
\item $C$ is a third position acting on the relation,
\item $d \in [0,\infty)$ is a distance parameter.
\end{itemize}
\end{definition}

\section{Coupling}

Define coupling:
\[
k = \frac{1}{d + \epsilon}.
\]

\begin{itemize}
\item $d \to \infty \Rightarrow k \to 0$ (negligible influence)
\item $d \to 0 \Rightarrow k \to \infty$ (dominant influence)
\end{itemize}

\section{Regimes of the Third Term}

\begin{definition}[Modes of C]
The third position $C$ operates in two regimes:
\begin{itemize}
\item \textbf{Noise}: stochastic perturbation
\item \textbf{Metastability}: structured persistence
\end{itemize}
\end{definition}

\begin{proposition}
Noise and metastability are not distinct ontologies but phase regimes of the same position.
\end{proposition}

\section{Transformation of Relations}

Let $R(A,B)$ denote the relation strength.

Define:
\[
R' = T(R, C, k)
\]

\begin{align*}
\text{Noise: } & R' = R + \eta(C)\,k \\
\text{Metastability: } & R' = R + \alpha C k
\end{align*}

\section{Role Permutation}

\begin{definition}[Permutation]
Let $\sigma$ be a permutation of $(A,B,C)$:
\[
(A,B,C) \mapsto \sigma(A,B,C)
\]
\end{definition}

\begin{proposition}
Roles are positional and not intrinsic.
\end{proposition}

\section{Dynamics}

The machine evolves via:
\[
(A,B,C)_{t+1} = \Phi(A,B,C,d)
\]

where $\Phi$ includes:
\begin{itemize}
\item transformation of relation,
\item update of positions,
\item optional permutation,
\item adjustment of distance.
\end{itemize}

\section{Emergence}

\begin{proposition}
The Parasite Machine generates:
\begin{itemize}
\item clusters,
\item regimes,
\item metastable entities,
\item transitions between order and noise,
\end{itemize}
without requiring conserved quantities.
\end{proposition}

\section{Interpretation}

The Parasite Machine demonstrates:
\begin{itemize}
\item structure emerges from interference,
\item entities are stabilized perturbations,
\item ontology is produced, not assumed.
\end{itemize}

\section{Relation to Serres}

The machine formalizes key insights from Michel Serres:
\begin{itemize}
\item the parasite as noise and mediator,
\item the transformation of relations through interference,
\item the instability of structure.
\end{itemize}

\section{Conclusion}

We have defined a minimal relational logic in which:
\begin{itemize}
\item no substance is assumed,
\item no conservation is required,
\item structure emerges from relational transformation.
\end{itemize}

\begin{quote}
A world may be constructed from relations acted upon by a third term whose regime determines whether it disrupts or stabilizes.
\end{quote}

\end{document}
